\documentclass[10pt,a4paper]{amsart}
\usepackage[margin=19mm]{geometry}
\usepackage{amsmath,amssymb,mathtools}
\usepackage[scr=rsfs,cal=euler]{mathalfa}
\usepackage{xfrac}
\usepackage[unicode,hidelinks,bookmarks=false]{hyperref}
\newcommand{\ie}{i.e.\ }
\newcommand{\cf}{cf.\ }
\newcommand{\resp}{respectively\ }
\newcommand{\Subsection}{\subsection}
\providecommand{\nobreakdash}{\nobreak\hbox{-}}
\setlength{\emergencystretch}{2em}
\multlinegap0pt
\usepackage{bm}
\usepackage{bbm}
\usepackage{dsfont}
\usepackage{xspace}
\usepackage{cancel}
\usepackage{mathtools}
\usepackage{cleveref}
\usepackage{amsthm}
\makeatletter
\@ifundefined{assumption}{\newtheorem{assumption}{Assumption}}{}
\@ifundefined{theorem}{\newtheorem{theorem}{Theorem}}{}
\makeatother
\newcommand{\abs}[1]{\left| #1\right|}
\newcommand{\bbE}{\mathbb{E}}
\newcommand{\bbP}{\mathbb{P}}
\newcommand{\calH}{{\mathcal{H}}}
\newcommand{\lbra}[1]{\left[#1\right]}
\newcommand{\mbra}[1]{\left\{#1\right\}}
\newcommand{\sbra}[1]{\left(#1\right)}
\title[VGB for Masked Diffusion Model: Efficient Test-time Scaling ]{VGB for Masked Diffusion Model: Efficient Test-time Scaling for Reward Satisfaction and Sample Editing}
\author{Kijung Jeon, Thuy-Duong Vuong, Molei Tao}
\date{Source: arXiv:2606.28301v1 (CC BY 4.0)}
\begin{document}
\maketitle

\noindent\emph{Source and licence.} VGB for Masked Diffusion Model: Efficient Test-time Scaling for Reward Satisfaction and Sample Editing. Version arXiv:2606.28301v1 (CC BY 4.0), distributed under the Creative Commons Attribution 4.0 International licence, as recorded in the arXiv licence metadata for this exact version and shown on the abstract page https://arxiv.org/abs/2606.28301v1. Full rights notice: licenses/arxiv-2606.28301-CC-BY-4.0.txt.\par\smallskip

\noindent\textit{Editorial scope.} Counted unit: the Theorem whose source label is \texttt{thm:robust\_flow\_cancelled\_hitting} in the article (source0.tex lines 4268--4300, source label \texttt{thm:robust\_flow\_cancelled\_hitting}), delivered together with the article's own complete original proof of it (source0.tex lines 4302--4434, one \texttt{proof} environment of 133 lines). No mathematics is written, repaired or re-derived: statement and proof are the article's own text line for line. Required context retained uncounted: ass:directional\_near\_balance (lines 4117--4127). Every definition, display and label of those lines is retained unchanged. Omitted from this package and not redistributed: 1--4116, 4128--4267, 4301--4301, 4435--5450. Nothing inside a retained excerpt is omitted or altered: every retained body is delivered exactly as the article's lines are written, so no omission receipt of any kind accompanies this package. Citations: the retained excerpts carry no citation command at all, so no bibliography entry of the article is transcribed and no reference is redistributed. Reference adaptation: none was needed and none was made; every reference inside the delivered unit resolves inside it, so no unresolved reference or citation is produced. Printed slips kept literal: no printed word, symbol, hypothesis or number of the counted statement or of its proof was altered. Layout and class: the article's own class is \texttt{article} with packages including natbib, neurips\_2026, inputenc, fontenc, pdfrender, url, booktabs, amsfonts, nicefrac, microtype, makecell, subfiles, subcaption, comment; the wrapper is the lane's pinned \texttt{amsart} editorial wrapper at 10pt on A4, which restates the theorem-like environments the retained text needs and adds the article's own macro definitions that the retained text needs (\texttt{\textbackslash abs}, \texttt{\textbackslash bbE}, \texttt{\textbackslash bbP}, \texttt{\textbackslash calH}, \texttt{\textbackslash lbra}, \texttt{\textbackslash mbra}, \texttt{\textbackslash sbra}), with extra wrapper packages (bm, bbm, dsfont, xspace, cancel, mathtools, cleveref). The delivered numbering is the unit's own. Assets: no retained span contains a figure environment, a graphics-inclusion command, a PDF-page inclusion, a table or a TikZ picture; no image file is copied into this package, the package declares no asset dependency and no image file is redistributed with it. Licence: version arXiv:2606.28301v1 of this article is distributed under the Creative Commons Attribution 4.0 International licence, as recorded in the arXiv licence metadata for this exact version and shown on the abstract page https://arxiv.org/abs/2606.28301v1. A full rights notice is delivered at licenses/arxiv-2606.28301-CC-BY-4.0.txt. Characters: every retained excerpt is pure ASCII, so the delivered engine, XeLaTeX with its default Latin Modern text font, reports no missing character.
\begin{assumption}[Uniform directional near-balance]
\label{ass:directional_near_balance}
There exists $\varepsilon_n\in[0,1)$ such that for every reachable interior state $z$,
\begin{equation*}
    \abs{p_{\downarrow}(z)-p_{\uparrow}(z)}\le \varepsilon_n.
\end{equation*}
Equivalently, since $p_{\downarrow}(z)+p_{\uparrow}(z)=1$,
\begin{equation*}
    \frac{1-\varepsilon_n}{2}\le p_{\downarrow}(z),p_{\uparrow}(z)\le \frac{1+\varepsilon_n}{2}.
\end{equation*}
\end{assumption}

\begin{theorem}[First-leaf hitting for near-maximal flow cancellation]
\label{thm:robust_flow_cancelled_hitting}
Assume \cref{ass:directional_near_balance}. Run the flow-cancelled momentum
chain $\widetilde K_{\mathrm{BAL}}^{\mathrm{fc},\chi}$ from the initial
lifted state $(\varnothing,\downarrow)$, where $\chi\in[0,1]$. Then
\begin{equation*}
    \bbE\lbra{
        T_n^{\mathrm{BAL,Mo}}
        \mid
        \widetilde Z_0=(\varnothing,\downarrow)
    }
    \le
    \sbra{
        \frac{4n}{1-\varepsilon_n}+1
    }
    \sbra{
        \frac{2-\chi(1-\varepsilon_n)}
             {1-\varepsilon_n}
    }^n .
\end{equation*}
In particular, if $\varepsilon_n\le \min\mbra{c/n,1/2}$ and
$1-\chi\le c_\chi/n$ for fixed constants $c,c_\chi>0$, then
\begin{equation*}
    \bbE\lbra{
        T_n^{\mathrm{BAL,Mo}}
        \mid
        \widetilde Z_0=(\varnothing,\downarrow)
    }
    =
    O_{c,c_\chi}(n).
\end{equation*}
For $\chi=1$, this recovers the maximally flow-cancelled bound.
\end{theorem}

\begin{proof}
Let $q_{\varepsilon}:=(1-\varepsilon_n)/2$. Under
\cref{ass:directional_near_balance}, at every interior state $z$,
$p_{\downarrow}(z),p_{\uparrow}(z)\ge q_{\varepsilon}$. From
$(z,\downarrow)$, the flow-cancelled kernel has downward move probability
$p_{\downarrow}(z)$, switch-up probability
$p_{\uparrow}(z)-\chi m(z)$, and same-copy holding probability
$\chi m(z)$, where
$m(z):=\min\{p_{\downarrow}(z),p_{\uparrow}(z)\}$. Thus, while in downward
mode, the probability of a downward depth increment is at least
$q_{\varepsilon}$.

We decompose the trajectory into downward attempts. A downward attempt starts
in downward mode. It succeeds if the chain reaches depth $n$ while remaining
in downward mode, and it fails if the chain switches from downward mode to
upward mode before reaching depth $n$.

Same-copy holding steps do not change either the depth or the momentum label.
For the success-probability calculation, condition them out and look only at
decisive steps, namely steps on which either a downward move or an upward switch
occurs. Given the current lifted state $(z,\downarrow)$,
\begin{align*}
    \bbP(\text{downward move}\mid \text{decisive},z,\downarrow)
    &=
    \frac{p_{\downarrow}(z)}
    {p_{\downarrow}(z)+p_{\uparrow}(z)-\chi m(z)}\\
    &=
    \frac{p_{\downarrow}(z)}
    {1-\chi m(z)}
    \ge
    \frac{q_\varepsilon}
    {1-\chi q_\varepsilon}
    :=
    \alpha_{\chi,\varepsilon}
    =
    \frac{1-\varepsilon_n}{2-\chi(1-\varepsilon_n)}.
\end{align*}
To succeed from depth $k$, an attempt only needs the next $n-k$ decisive
outcomes to be downward moves. Therefore, with
$s_{\chi,\varepsilon,n}:=\alpha_{\chi,\varepsilon}^{\,n}$,
\begin{equation*}
    \bbP(\text{attempt succeeds from depth }k)
    \ge
    \alpha_{\chi,\varepsilon}^{\,n-k}
    \ge
    \alpha_{\chi,\varepsilon}^{\,n}.
\end{equation*}
Thus every downward attempt succeeds with probability at least
$s_{\chi,\varepsilon,n}$, uniformly over the past.

Next, we bound the expected duration of one attempt cycle, counting Markov
transitions including same-copy holding steps. During the downward part, every
transition in downward mode increases the depth with probability at least
$q_\varepsilon$. To upper-bound the time cost, ignore early switch failures
and wait until $n$ downward depth increments have occurred. If $G_r$ is the
waiting time for the $r$-th such increment after the $(r-1)$-st one, then
$\bbP(G_r>t)\le (1-q_\varepsilon)^t$ and
$\bbE[G_r]\le 1/q_\varepsilon$. Hence the expected duration of the downward
part is at most $n/q_\varepsilon$.

If the attempt fails, the chain has switched to upward mode before reaching
depth $n$. We upper-bound the cost of preparing the next attempt by a
conservative reset: move upward until the root and then take the boundary
switch back to downward mode. Returning to the root requires at most $n$
upward depth decrements, each occurring with probability at least
$q_\varepsilon$. Thus the reset duration is at most $n/q_\varepsilon+1$
in expectation, and one complete attempt cycle has expected duration at most
$L_\varepsilon:=2n/q_\varepsilon+1$.

Let $N_{\mathrm{att}}$ be the number of attempts until the first successful
one. Since each attempt succeeds with conditional probability at least
$s_{\chi,\varepsilon,n}$, we have
$\bbP(N_{\mathrm{att}}>m)\le(1-s_{\chi,\varepsilon,n})^m$ and therefore
$\bbE[N_{\mathrm{att}}]\le 1/s_{\chi,\varepsilon,n}$.

Let $C_j$ denote the duration of the $j$-th attempt cycle. The cycle-length
bound above is uniform over the past, so
$\bbE[C_j\mid \calH_{j-1}]\le L_\varepsilon$ whenever the $j$-th cycle
starts. Since
$T_n^{\mathrm{BAL,Mo}}\le
\sum_{j\ge1}\mathbf 1\{N_{\mathrm{att}}\ge j\}C_j$,
we obtain
\begin{equation*}
    \bbE[T_n^{\mathrm{BAL,Mo}}]
    \le
    L_\varepsilon\sum_{j\ge1}\bbP(N_{\mathrm{att}}\ge j)
    =
    L_\varepsilon\bbE[N_{\mathrm{att}}]
    \le
    \left(\frac{2n}{q_\varepsilon}+1\right)
    \alpha_{\chi,\varepsilon}^{-n}.
\end{equation*}
Substituting $q_\varepsilon=(1-\varepsilon_n)/2$ and
$\alpha_{\chi,\varepsilon}=(1-\varepsilon_n)/(2-\chi(1-\varepsilon_n))$
gives
\begin{equation*}
    \bbE[T_n^{\mathrm{BAL,Mo}}]
    \le
    \sbra{
        \frac{4n}{1-\varepsilon_n}+1
    }
    \sbra{
        \frac{2-\chi(1-\varepsilon_n)}
             {1-\varepsilon_n}
    }^n .
\end{equation*}

Finally, suppose $\varepsilon_n\le \min\mbra{c/n,1/2}$ and
$1-\chi\le c_\chi/n$. Then
\begin{align*}
    \frac{2-\chi(1-\varepsilon_n)}
         {1-\varepsilon_n}
    &=
    1+
    \frac{
        1-\chi+\varepsilon_n(1+\chi)
    }{
        1-\varepsilon_n
    }=
    1+O_{c,c_\chi}\!\left(\frac1n\right).
\end{align*}
Therefore its $n$-th power is $O_{c,c_\chi}(1)$, while
$(4n)/(1-\varepsilon_n)+1=O_c(n)$. Hence
\begin{equation*}
    \bbE\lbra{
        T_n^{\mathrm{BAL,Mo}}
        \mid
        \widetilde Z_0=(\varnothing,\downarrow)
    }
    =
    O_{c,c_\chi}(n).
\end{equation*}
\end{proof}

\end{document}
